%% notes-tex/027-metabolic-flux-module/027-metabolic-flux-module.tex
%% [[experiments.027-betaxanthin-metabolic]]
%% https://github.com/Mjvolk3/torchcell/tree/main/notes-tex/027-metabolic-flux-module/027-metabolic-flux-module.tex
%%
%% The consolidated account of the enzyme-constrained thermodynamic flux module: what it
%% computes, which parameter tables it consumes, which of those are wired in, and what has
%% been measured. Assembled from the 026 and 027 dendron notes plus the module source.
%%
%% Build:  make            -> 027-metabolic-flux-module.pdf
%%         make check      -> formatting / provenance / style gate
%%         make plots      -> re-convert the matplotlib panels from ASSET_IMAGES_DIR
%%
%% Every number is produced by a committed script under experiments/026-metabolism-flux/
%% or experiments/027-betaxanthin-metabolic/ and read back from their results/ trees.

\documentclass[10pt,a4paper]{article}
\usepackage[draft]{tcdoc}
\usepackage{float}
\usepackage{amsmath}
\usepackage{amssymb}

\emergencystretch=3.5em

\title{\vspace{-1.5em}The Enzyme-Constrained Thermodynamic\\Flux Module\\
  {\large What it computes, what it is fed,\\and what has been measured}}
\author{Michael Volk}
\date{\today}

\begin{document}
\maketitle
\thispagestyle{empty}

\begin{abstract}
\noindent
The flux module is a differentiable layer that maps a perturbed gene representation to a
flux vector over yeast-GEM 9.0.2, subject to stoichiometry, an enzyme-capacity bound and a
thermodynamic sign constraint. It is neither a Type~I nor a Type~II virtual instrument as
those were defined: its output is a latent physical state rather than a measured phenotype,
and the phenotypes are read off that state.

Its distinguishing property is that three of its inputs are meant to be tabulated rather
than learned: turnover numbers, Michaelis constants and formation energies, each
hash-pinned and attributable to the weights or model version that produced it. That is what
would make the layer a reconstruction of a metabolic network rather than a parameterization
fitted to one organism.

\textbf{None of the three tables built in this work is yet consumed by any measured arm.}
The predicted $k_{cat}$ tables are not plumbed into the capacity bound, which still resolves
4.0 percent measured values from the Open Enzyme Database and defaults the rest; the
recomputed formation energies are not plumbed into the potential, which still reads
yeast-GEM's shipped file; and $K_M$ has no consumer at all. The module diagram draws those
edges as solid and the code does not have them. What exists is an audit of what the inputs
could be, built reproducibly enough to apply to another organism, and every result here
should be read as testing the module's structure rather than its parameters.

This document consolidates what was scattered across the 026 and 027 notes: the
formulation as implemented, the three parameter tables and the sharp distinction between
those wired in and those merely built, the ablation registry, and every measured number.
Two corrections to earlier claims are recorded rather than quietly folded in. The
ChemAxon licensing blocker does not apply to the betaxanthin pathway, whose intermediates
are already in the compound cache. And the reaction-coverage shortfall is not a
mapping gap: an audit that recovers every identifier route available gains 0.58
percentage points, because 375 metabolites match a cache record that carries no structure
and are counted as resolved while returning no energy.

The training runs that test whether the module helps are on Delta as this is written, and
the preliminary section reports only what has actually been observed.
\end{abstract}

\vspace{0.5em}
\tccontents
\clearpage

%% notes-tex/027-metabolic-flux-module/sections/1-what-it-is.tex
%% [[experiments.027-betaxanthin-metabolic]]
%% https://github.com/Mjvolk3/torchcell/tree/main/notes-tex/027-metabolic-flux-module/sections/1-what-it-is.tex

\section{What the module computes\secstatus{todo}}
\label{sec:what}

\subsection{Where it sits\secstatus{todo}}

The Equivariant Cell Graph Transformer separates its instruments by what they map between.
A \textbf{Type~I} instrument is representation to representation: the perturbation
operators $T_\psi^{\mathrm{del}}$, $T_\psi^{\mathrm{OE}}$, $T_\psi^{\mathrm{KD}}$ act
inside the gene embedding space and hand on $H_{\mathrm{pert}}$. A \textbf{Type~II}
instrument is representation to output: a readout that emits a measured phenotype.

The flux module is neither, and the reason is worth stating precisely rather than settling
by convention. By the letter of the taxonomy it is Type~I, since its output feeds further
computation rather than a loss. But every existing Type~I instrument acts \emph{within} the
gene index space at dimension $d$ and preserves equivariance over genes, whereas this one
leaves gene index space entirely for reaction space $\mathbb{R}^{4131}$ under a fixed,
externally supplied stoichiometry, and its parameters are tabulated physical constants
rather than learned weights. It is a physically identified decoder, and its output
$\nu$ is a latent physical state rather than a phenotype: the phenotypes are read
\emph{off} it.

\subsection{A program cannot learn, which is the point\secstatus{todo}}

The constraint-based tradition solves a linear program per condition. A linear program has
no use for a phenotype measurement, so it cannot be improved by one. This layer inverts
that: every constraint is either enforced exactly by the parameterization or relaxed to a
smooth penalty, so the flux is \emph{conditioned} on roughly $10^7$ fitness records, 4,735
betaxanthin measurements and 4,678 amino-acid profiles. \textbf{Nothing is enforced by a
binary variable}, which is what keeps the whole object differentiable and lets a batch of
perturbations cost the same forward pass as a single one.

\subsection{The relaxation, term by term\secstatus{todo}}

\begin{table}[H]
  \centering
  \footnotesize
  \caption[The relaxation]{Each constraint of an enzyme-constrained thermodynamic model,
  and how this layer realizes it. ``Exact'' means the parameterization cannot violate it;
  ``soft'' means a penalty that the optimizer trades against the data term.}
  \label{tab:relaxation}
  %% SOURCE: torchcell/metabolism/flux_layer.py, module docstring table and forward().
  \begin{tabular}{@{}p{0.26\textwidth}p{0.34\textwidth}p{0.32\textwidth}@{}}
    \hline
    Constraint & Conventional form & Here \\
    \hline
    box and directionality & $lb_j \le \nu_j \le ub_j$ & \textbf{exact}, via sigmoid \\
    enzyme capacity & $\nu_j \le k_{cat} E$ & \textbf{exact}, folded into the box \\
    mass balance & $S\nu = 0$ & soft, scale-relative \\
    second law & $\Delta_r G_j \nu_j < 0$ plus a binary & soft hinge, no binary \\
    Gibbs dissipation & $g_{\mathrm{diss}} \le g_{\mathrm{lim}}$ & soft hinge \\
    protein budget & $\sum_g MW_g E_g \le P_{\mathrm{avail}}$ & soft hinge \\
    \hline
  \end{tabular}
\end{table}

\paragraph{From representation to gene availability.} Each gene's embedding passes through
a learned gate, $\gamma_g = \sigma(w^\top h_g)$, so a gene is not merely present or absent.
A deletion pins $\gamma_g = 0$; everything else is free to take an intermediate value,
which is what leaves room for dosage, hypomorphic alleles and overexpression rather than
forcing the layer to treat every perturbation as a knockout.

\paragraph{Folding up the gene-protein-reaction rule.} Availability is folded to the
reaction through the GPR: a \textbf{softmin} over the subunits of a complex, since a
complex is limited by its scarcest member, and a \textbf{sum} over isozymes, since they
substitute for one another. The softmin temperature is a real parameter, not a detail: too
soft and a complex behaves like a mean over its subunits, which loses exactly the
limiting-subunit behavior the softmin was chosen for.

\paragraph{The box, which is where capacity lives.} Flux is parameterized as
\begin{equation}
  \nu_j = lb_j + (ub_j - lb_j)\,\sigma(z_j),
  \qquad
  |\nu_j| \le k_{cat,j}\,E_{g(j)},
\end{equation}
so the bounds and the capacity bound hold by construction and cannot be violated no matter
what the network does. Turnover enters in units of $\mathrm{h}^{-1}$
($k_{cat}\,[\mathrm{s}^{-1}] \times 3600$), which is what fixes the units of the enzyme
abundance $E$. The exactness budget is spent here rather than on $S\nu = 0$, and that is a
deliberate choice: bounds are exact and mass balance is soft, not the reverse.

\paragraph{The second law, which is the load-bearing relaxation.}
\begin{equation}
  C_{\mathrm{th}} = \frac{1}{|\mathcal{J}|}\sum_{j\in\mathcal{J}}
    \operatorname{relu}\!\big(\nu_j\,\Delta_r G_j + \epsilon\big),
\end{equation}
with $\mathcal{J}$ excluding exchange reactions, the biomass reaction and water transport,
which cross the system boundary or are lumped and so carry no meaningful reaction energy.
The $\epsilon$ is not decoration. The conventional big-M pair admits
$\nu_j = \Delta_r G_j = 0$ under either value of its binary, so the strict inequality is
never actually realized, and a smooth version has to name its tolerance rather than inherit
a fiction.

\textbf{Loop-freedom comes free, and that is the trick.} Because $\Delta_r G$ is a
difference of potentials it sums to zero around any cycle, so no cycle can run downhill at
every step. Requiring $\nu_j \Delta_r G_j \le 0$ therefore forbids internal loops with no
integer variable at all, which is precisely the property that makes loopless flux balance
analysis a mixed-integer program.

\subsection{Free against anchored thermodynamics\secstatus{todo}}

The in-cell reaction energy decomposes into three terms,
\begin{equation}
  \Delta_r G_j
  = \underbrace{\Delta_r G'^{\circ}_j}_{\text{tabulated}}
  + \underbrace{RT\!\!\sum_{i \neq \mathrm{H}^+}\!\! S_{ij} \ln c_i}_{\text{concentration}}
  + \underbrace{\Delta_r G^{t}_j}_{\text{transport}} .
\end{equation}

The distinction between the two thermodynamic arms is exactly which of these is supplied.
In the \textbf{free} arm the potential is learned, so the layer gets loop-freedom and
nothing else; the constraint is structural. In the \textbf{anchored} arm the standard term
is read from the tabulated formation energies, so the sign of every reaction is pinned by
measured chemistry rather than by whatever the optimizer finds convenient. That difference
is the experiment: structural loop-freedom is a real constraint on its own, and the
question is whether the chemistry adds anything to it.

%% notes-tex/027-metabolic-flux-module/sections/2-parameters.tex
%% [[experiments.026-metabolism-flux.kinetic-and-thermodynamic-parameter-datasets]]
%% https://github.com/Mjvolk3/torchcell/tree/main/notes-tex/027-metabolic-flux-module/sections/2-parameters.tex

\section{The three parameter tables\secstatus{todo}}
\label{sec:parameters}

\subsection{What was there before, and what is there now\secstatus{todo}}

The module needs three tables and had almost none of them. Measured $k_{cat}$ covered 148
of 3{,}728 catalytic units, 4.0 percent, with the other 96 percent running on a single flat
default of 10.3\,s$^{-1}$. $K_M$ had no table at all: the parameter existed as an enum
member and a protocol flag with nothing behind either. The thermodynamics came from a
shipped file of point values with no covariance, usable only for this organism.

Six kinetic predictors were mirrored into
\texttt{\$DATA\_ROOT/models/kinetics/<name>/} with per-file \texttt{sha256}, retrieval
commands and git commits, and all six now run over the enzyme-substrate pairs of yeast-GEM
9.0.2, with protein sequences taken from the genome object rather than from the table each
model bundles.

\begin{table}[H]
  \centering
  \footnotesize
  \caption[The built tables]{Every parameter table now built, with the exponentiation base
  established by validating each predictor against its own authors' published output.
  Coverage is of the GEM's 3{,}728 catalytic units.}
  \label{tab:tables}
  %% SOURCE: experiments/026-metabolism-flux/results/kinetics_<predictor>_<param>_summary.json
  \begin{tabular}{@{}llrrrl@{}}
    \hline
    Predictor & Parameter & Rows & Units & Coverage & Median \\
    \hline
    DLKcat & $k_{cat}$ & 7{,}351 & 3{,}552 & 95.3\% & 4.012 s$^{-1}$ \\
    UniKP & $k_{cat}$ & 7{,}351 & 3{,}552 & 95.3\% & 4.050 s$^{-1}$ \\
    EITLEM & $k_{cat}$ & 7{,}456 & 3{,}552 & 95.3\% & 3.981 s$^{-1}$ \\
    TurNuP & $k_{cat}$ & 5{,}143 & 2{,}056 & 55.2\% & 10.819 s$^{-1}$ \\
    DeepEnzyme & $k_{cat}$ & 7{,}300 & 3{,}546 & 95.1\% & 4.167 s$^{-1}$ \\
    \hline
    UniKP & $K_M$ & 7{,}351 & 3{,}552 & 95.3\% & 0.0796 mM \\
    EITLEM & $K_M$ & 7{,}456 & 3{,}552 & 95.3\% & 0.0793 mM \\
    Boost\_KM & $K_M$ & 7{,}157 & 3{,}468 & 93.0\% & 0.1410 mM \\
    \hline
  \end{tabular}
\end{table}

\textbf{Two of these models exponentiate base 2 and four base 10.} DLKcat and DeepEnzyme
are $\log_2$; UniKP, EITLEM, TurNuP and Boost\_KM are $\log_{10}$. Reading either wrong
yields turnover numbers that still look entirely credible, which is why each was validated
against something its own authors published before its values were believed.

\subsection{Which tables are actually wired in\secstatus{todo}}

This distinction decides how any result from the module should be read, and it is easy to
lose.

\begin{table}[H]
  \centering
  \footnotesize
  \caption[Wired against built]{What the layer consumes today. A table that is built but
  not consumed contributes nothing to a measured result and must not be described as if it
  did.}
  \label{tab:wired}
  %% SOURCE: torchcell/metabolism/flux_layer.py __init__ and build_flux_layer in
  %% experiments/026-metabolism-flux/scripts/train_flux.py.
  \begin{tabular}{@{}p{0.16\textwidth}p{0.30\textwidth}p{0.46\textwidth}@{}}
    \hline
    Table & Status & How it enters, or why it does not \\
    \hline
    $k_{cat}$ & \textbf{wired}, but from the Open Enzyme Database, not the predictors &
      the capacity bound $|\nu_j| \le k_{cat,j} E_{g(j)}$. \texttt{resolve\_kcat\_table}
      reads the OED mirror, so the run uses 4.0 percent measured values and a default for
      the rest. The predicted tables are not yet plumbed to this call \\
    $\Delta_f G'^{\circ}$ & \textbf{wired} in the anchored arm &
      the standard term of $\Delta_r G_j$, hence the second-law hinge and the dissipation
      bound \\
    $K_M$ & \textbf{built, not wired} &
      the saturation term is drawn as inactive in the module diagram and is inactive in the
      code. $K_M$ contributes nothing to any result reported here \\
    \hline
  \end{tabular}
\end{table}

The consequence is worth stating plainly: the current runs test the module's
\emph{structure and thermodynamics}, not its predicted kinetics. Wiring the predictor
registry into \texttt{resolve\_kcat\_table} is the next change, and it is the one that
turns the five kcat tables from an audit into an input.

\subsection{Why five predictors and not one\secstatus{todo}}

Re-keying the comparison from the enzyme-substrate pair to the \textbf{gene}, which is the
unit the capacity bound reads, gives the number that decides how a predicted table should
be used.

\begin{table}[H]
  \centering
  \footnotesize
  \caption[Per-gene disagreement]{Across-predictor spread per gene against the dynamic
  range of the quantity itself, in decades, over genes every predictor covers.}
  \label{tab:pergene}
  %% SOURCE: experiments/026-metabolism-flux/results/per_gene_agreement_{k_cat,K_M}.json
  \begin{tabular}{@{}lrrrr@{}}
    \hline
    Quantity & Genes & Median spread & Consensus p10--p90 & Ratio \\
    \hline
    $k_{cat}$ & 980 & \textbf{1.23} & 0.94 & \textbf{1.32} \\
    $K_M$ & 1{,}150 & 0.70 & 1.34 & 0.52 \\
    \hline
  \end{tabular}
\end{table}

\textbf{For $k_{cat}$ the disagreement between predictors exceeds the quantity being
predicted.} The median gene moves 1.23 decades between the highest and lowest predictor, a
factor of 17, and 67 percent of genes move by more than a full decade, against a
10th-to-90th-percentile spread across genes of 0.94 decades. Choosing a different published
predictor moves one gene's turnover number further than real enzymes differ from one
another. Gene-level Spearman between $k_{cat}$ predictors runs 0.04 to 0.48.

\textbf{For $K_M$ it does not.} The same construction gives 0.70 against 1.34, and
Boost\_KM agrees with UniKP at a gene-level Spearman of 0.75, higher than any pair of
$k_{cat}$ predictors reaches. Both are statements about agreement rather than accuracy,
since almost none of these pairs has a measured value.

That is why the registry in \texttt{torchcell/datasets/kinetics/} exists and why a config
string selects the predictor: predictor choice is a hyperparameter with a measured effect
size, not a settled detail.

%% notes-tex/027-metabolic-flux-module/sections/3-coverage.tex
%% [[experiments.026-metabolism-flux.scripts.audit_unresolved_metabolites]]
%% https://github.com/Mjvolk3/torchcell/tree/main/notes-tex/027-metabolic-flux-module/sections/3-coverage.tex

\section{Thermodynamic coverage, and two corrections\secstatus{todo}}
\label{sec:coverage}

\subsection{The recomputation\secstatus{todo}}

Formation energies were recomputed from eQuilibrator's component contribution, which
returns a covariance and applies to any organism because it works from structure rather
than from a curated per-organism file. Summing the resulting table reproduces
eQuilibrator's own reaction call to $9.1\times10^{-13}$\,kJ\,mol$^{-1}$, so the table is
arithmetically correct. On the chemical reactions both tables cover it agrees with the
shipped file to a median of 10.11\,kJ\,mol$^{-1}$.

The gain is 785 reactions with a real standard error, median 2.30\,kJ\,mol$^{-1}$, and the
rank-590 covariance factor the layer had recorded as unavailable. The cost is reaction
coverage: 42.5 percent against the shipped file's 77.7.

\subsection{Correction: the shortfall is not a mapping gap\secstatus{todo}}

The obvious reading of that 42.5 percent is that identifiers are failing to resolve, and
that better mapping would close it. An audit that implements every available recovery route
measures otherwise.

\textbf{There are two loss channels, and the summary reported only one.} The 332
``unresolved'' metabolites are the accession-match failures. A further \textbf{375
metabolites match a cache record that carries no structure and no group vector}, so
\texttt{standard\_dg\_formation} returns nothing while the tier histogram counts them as
resolved. Total unpriced is \textbf{707}, not 332. The baseline is 2{,}099 of 2{,}806
metabolites, 74.80 percent, and 1{,}755 of 4{,}131 reactions, 42.48 percent.

\begin{table}[H]
  \centering
  \footnotesize
  \caption[Measured recovery]{Every identifier route available, applied in turn. Reaction
  coverage is the binding number because a reaction needs every participant priced.}
  \label{tab:recovery}
  %% SOURCE: experiments/026-metabolism-flux/results/unresolved_metabolites.json
  \begin{tabular}{@{}lrrr@{}}
    \hline
    Step & Metabolites priced & Reactions & $\Delta$ reactions \\
    \hline
    baseline & 2{,}099 (74.80\%) & 1{,}755 (42.48\%) & -- \\
    extra identifier namespaces & 2{,}102 & 1{,}764 (42.70\%) & $+9$ \\
    MetaNetX cross-reference table & 2{,}113 & 1{,}770 (42.85\%) & $+6$ \\
    InChIKey structure lookup & 2{,}127 (75.80\%) & 1{,}779 (43.06\%) & $+9$ \\
    \hline
  \end{tabular}
\end{table}

\textbf{Total recovery is 28 metabolites and 24 reactions, 0.58 percentage points.} The
extra namespaces are worth exactly one chemical, hydrogen sulfide. The census explains why:
the model carries chebi on 2{,}402 metabolites and metanetx on 2{,}251, but biocyc on 29,
seed on 16, reactome on 2, and no hmdb, lipidmaps or swisslipid at all. Of the 332
unmatched, \textbf{181 of the 192 distinct chemicals carry no chemical annotation
whatsoever}, and the class breakdown is overwhelmingly lipid: phosphatidylinositol 135,
cardiolipin 62, triacylglycerol 34, phosphatidate 25.

Of the 679 still unpriced, 249 (36.7 percent) are a real ceiling: 185 carry a placeholder
formula with an R, X or $*$ so no definite structure exists, and 64 are known to
eQuilibrator but have no group decomposition. The other 430 (63.3 percent) have a definite
formula and no usable structure anywhere, which is a structure-curation gap rather than a
limit of the method.

\subsection{A correctness bug in the shipped build, found by the same guard\secstatus{todo}}

The InChIKey route initially reported a gain six times larger, and it was wrong. MetaNetX
writes \texttt{[*]} dummy atoms for polymers and conjugates; RDKit returns an \emph{empty}
InChIKey for these rather than failing; an empty key reaches the cache as \texttt{LIKE
'\%'} and matches the entire compound table. \texttt{Ala-tRNA(Ala)} was being priced as
water.

The guard that fixes it, comparing heavy-atom composition against the GEM formula with
hydrogen excluded because protonation state legitimately differs, then finds the same class
of error in the build already in use: \textbf{31 metabolites currently carry a formation
energy from a compound whose heavy-atom composition disagrees with the GEM formula}, and 14
of the 1{,}755 covered reactions touch one. Every dolichol species is matched to a C25
pentaprenyl against the GEM's C20 tetraprenyl; \texttt{glycogen} is matched to a C24
tetramer against a C6 monomer; \texttt{protein}, \texttt{carbohydrate} and \texttt{sterols}
are matched to concrete molecules.

This does not raise coverage and it is worth more than coverage: it removes 31 wrong
numbers. It is not yet applied to \texttt{resolve\_compound}, and it should be.

\subsection{Correction: the ChemAxon blocker does not apply here\secstatus{todo}}

An earlier account of this work stated that pricing a novel compound requires ChemAxon's
\texttt{cxcalc} for pKa estimation, that no license is available, and that the betaxanthin
intermediates therefore cannot be priced. The first two clauses are true and the conclusion
is false: \textbf{all of the pathway's intermediates are already in the compound cache}, so
nothing about betaxanthin needs a new compound created.

For a genuinely novel compound the constraint has also loosened. Upstream
\texttt{equilibrator-assets} merged support on 2026-08-15 for creating compounds without
ChemAxon when pKa values are supplied, via \texttt{pkas\_provided\_externally} and
\texttt{specified\_pkas}. That is master-only; the pinned release 0.6.0 does not contain it.
What the interface needs is an unannotated list of macroscopic pKa floats, which is a
weaker requirement than the site-annotated micro-pKa the harder tools produce.

%% notes-tex/027-metabolic-flux-module/sections/4-arms.tex
%% [[experiments.027-betaxanthin-metabolic]]
%% https://github.com/Mjvolk3/torchcell/tree/main/notes-tex/027-metabolic-flux-module/sections/4-arms.tex

\section{The ablation registry\secstatus{todo}}
\label{sec:arms}

\subsection{What each arm turns on\secstatus{todo}}

The registry is one dictionary, \texttt{ARMS} in
\texttt{experiments/026-metabolism-flux/scripts/train\_flux.py}, and it is the single
source of truth for what the layer does in any run. An experiment names an arm; it never
redefines one. That is what keeps a difference between two arms attributable to the module
rather than to a hyperparameter that moved alongside it.

\begin{table}[H]
  \centering
  \footnotesize
  \caption[The arms]{The registry as implemented. Each row is a strict addition to the one
  above it except \texttt{nullspace}, which changes the parameterization rather than the
  physics.}
  \label{tab:arms}
  %% SOURCE: experiments/026-metabolism-flux/scripts/train_flux.py, ARMS.
  \begin{tabular}{@{}lllll@{}}
    \hline
    Arm & Flux & Thermo & Capacity, budget, dissipation & Isolates \\
    \hline
    \texttt{pooled} & off & -- & -- & the no-module control \\
    \texttt{flux\_off} & on & \texttt{OFF} & off & routing through $S$ alone \\
    \texttt{flux\_free} & on & \texttt{FREE} & off & structural loop-freedom \\
    \texttt{flux\_anchored} & on & \texttt{ANCHORED} & on & \textbf{the treatment} \\
    \texttt{flux\_nullspace} & on & \texttt{ANCHORED} & on & the exactness budget \\
    \hline
  \end{tabular}
\end{table}

The ladder is deliberate. \texttt{flux\_off} routes the representation through the
stoichiometric network with nothing but mass balance and the GPR, so it asks whether the
network's topology alone carries signal. \texttt{flux\_free} adds a learned potential,
which buys loop-freedom and no chemistry. \texttt{flux\_anchored} replaces the learned
potential with tabulated formation energies and switches on the enzyme capacity bound, the
protein budget and the dissipation limit. \texttt{flux\_nullspace} keeps the anchored
physics but spends the exactness budget on $S\nu = 0$ instead of on the box, which is the
one arm that tests the parameterization choice rather than the model.

\subsection{The control that is missing, and why it matters\secstatus{todo}}

Every arm above changes the module's \emph{physics} together with its \emph{parameter
count}. So a gap between \texttt{flux\_anchored} and \texttt{pooled} is consistent with two
different explanations: the tabulated chemistry is informative, or the extra structure acts
as a capacity constraint that would help with any numbers in it.

The arm that separates those is \texttt{flux\_shuffled}: the full anchored module with
$k_{cat}$ \textbf{permuted across reactions}, holding architecture, parameter count and
optimization fixed and destroying only the correspondence between a reaction and its
turnover number. If \texttt{flux\_anchored} and \texttt{flux\_shuffled} land together, the
kinetic tables are decoration.

It is \textbf{not in the registry and not in the current run.} Adding it is a code change
that was not worth making untested, and it is the first item of the next round.

\subsection{The permutation nulls\secstatus{todo}}

Two arms exist only to be compared against: \texttt{null\_pooled} and
\texttt{null\_anchored} train on \textbf{permuted labels} while validating on real ones.
They matter because of what 026 measured. Against zero, its flux-versus-pooled gap of
$+0.0735$ is 7.7 standard errors. Against its own permuted-label gap of $+0.0302$, the
excess is $+0.0433$ at 1.9 standard errors. More than 40 percent of ``the flux module
helps'' survived destroying the labels.

So the decision rule is a difference of differences,
\begin{equation}
  \big(\text{flux} - \text{pooled}\big) - \big(\text{null flux} - \text{null pooled}\big),
\end{equation}
and a gap measured against zero is not evidence.

%% notes-tex/027-metabolic-flux-module/sections/5-measured.tex
%% [[experiments.026-metabolism-flux.enzyme-constrained-thermodynamic-flux-layer]]
%% https://github.com/Mjvolk3/torchcell/tree/main/notes-tex/027-metabolic-flux-module/sections/5-measured.tex

\section{What has been measured\secstatus{todo}}
\label{sec:measured}

\subsection{The 026 sweep, and why it decided nothing\secstatus{todo}}

\begin{table}[H]
  \centering
  \footnotesize
  \caption[026 arm results]{The 026 sweep on betaxanthin. The statistic is validation
  Pearson scored as a \emph{maximum over epochs}, which is an upward-biased order
  statistic, and the last column is the mean of the final five epochs, which is not.}
  \label{tab:026}
  %% SOURCE: experiments/026-metabolism-flux/results/flux_arms_summary.json
  \begin{tabular}{@{}lrrrrr@{}}
    \hline
    Arm & Seeds & Peak mean & s.e.m. & Peak epoch & Last-5 mean \\
    \hline
    \texttt{flux\_anchored} & 10 & 0.1151 & 0.0095 & 17 & 0.0057 \\
    \texttt{flux\_free} & 10 & 0.1020 & 0.0115 & 12 & 0.0212 \\
    \texttt{flux\_off} & 10 & 0.0837 & 0.0112 & 10 & 0.0044 \\
    \texttt{pooled} & 10 & 0.0416 & 0.0098 & 4 & $-0.0041$ \\
    \texttt{flux\_nullspace} & 10 & 0.0381 & 0.0114 & 0 & 0.0000 \\
    \hline
  \end{tabular}
\end{table}

The ordering looks like the ladder the design predicted. It is not evidence, for three
measured reasons.

\textbf{No arm clears its own null.} A label-permutation null over 24 runs has peak mean
0.0962, standard deviation 0.0441, and a 95th percentile of \textbf{0.1545}. No arm reaches
that percentile. The smallest $p$ across all 22 configurations tested is 0.24, and the
single highest peak anywhere in the sweep, 0.2110, belongs to a \emph{null} run.

\textbf{The last-5 column is the honest one.} Every arm's peak sits far above its own
end-of-training value, and three of five end negative. A maximum over 30 epochs of a
statistic whose analytic null width at $n_{\mathrm{val}} \approx 340$ is 0.0546 inflates
that width empirically to 0.0962, a factor of 1.8.

\textbf{The split moved with the seed.} 026 re-rolled the train/validation partition per
seed, so seed varied the nuisance and the signal together. Its paired arm gap has an
across-seed standard deviation of 0.0303 with real labels and 0.0722 with permuted ones,
both inflated by the partition moving underneath the contrast.

\subsection{The defect that outranks all three\secstatus{todo}}

Every betaxanthin run in this repository falls into one of two architecture families, and
the family accounts for nearly all of the variance between runs.

\begin{table}[H]
  \centering
  \footnotesize
  \caption[The two families]{Measured validation Pearson by backbone. The best betaxanthin
  result ever obtained is 0.43399, GilaHyper job 1344, config \texttt{gh\_metabolism\_002},
  seed 42, peak at epoch 39 of 80.}
  \label{tab:families}
  %% SOURCE: experiments/019-simb-multimodal/conf/gh_metabolism_00{0,2}.yaml and the
  %% 020 Optuna studies; see experiments/027-betaxanthin-metabolic/README.md.
  \begin{tabular}{@{}p{0.12\textwidth}p{0.56\textwidth}l@{}}
    \hline
    Family & Configuration & Measured \\
    \hline
    weak & 2 graphs, learnable embeddings, MSE, hidden 32, 4 heads & 0.04 to 0.16 \\
    strong & 9 graphs, \texttt{prot\_T5\_all}, no learnable embedding, CRPS, hidden 90,
      9 heads & 0.32 to \textbf{0.43} \\
    \hline
  \end{tabular}
\end{table}

\texttt{train\_flux.py} builds the \textbf{weak} family. So 026 compared a flux module
against a CGT running roughly five times below its own demonstrated ceiling, and its
\texttt{flux\_off} arm at 0.0837 is the correct control \emph{for that sweep} rather than a
CGT baseline. Nothing in 026 tested whether the module helps a CGT that works, which is
what the current run exists to do.

\subsection{The baselines the module has to beat\secstatus{todo}}

\begin{table}[H]
  \centering
  \footnotesize
  \caption[Baselines]{Every measured comparator, with the statistic named. FCL is
  Merzbacher's Flux Cone Learning, a three-class classifier rather than a regressor, scored
  on its own 639-gene test split.}
  \label{tab:baselines}
  %% SOURCE: experiments/020-cachera-betaxanthin/results/merzbacher_comparison_figures.json
  %% and merzbacher_head_to_head.json.
  \begin{tabular}{@{}p{0.30\textwidth}p{0.40\textwidth}l@{}}
    \hline
    Comparator & Statistic & Value \\
    \hline
    CGT, strong backbone & max val Pearson, $n_{\mathrm{val}} = 340$ & 0.43399 \\
    FCL \texttt{RandomForest\_Resampled} & gene-level accuracy (majority 0.674) & 0.700 \\
    FCL & Matthews correlation & 0.232 \\
    FCL & Pearson against the raw screen & 0.2435 \\
    FCL & Spearman against the raw screen & 0.0391 \\
    fully-connected / MLP baseline & \textbf{not measured; does not exist} & -- \\
    \hline
  \end{tabular}
\end{table}

There is no MLP baseline for betaxanthin anywhere in the repository. If one is wanted it
must be built, and no number should be quoted for it until it is.

\subsection{The noise floor, which sets the replicate count\secstatus{todo}}

Four independent measurements of the across-seed standard deviation on this endpoint:
0.0302 from hyperparameter-matched reruns, 0.0355 from ten fresh seeds of
\texttt{flux\_off}, 0.0299 from ten seeds of \texttt{flux\_anchored}, and 0.0337 from a
three-seed 019 arm. The analytic single-correlation width at $n_{\mathrm{val}} \approx 340$
is 0.0546.

Two consequences were computed rather than asserted. Selection inflation over 36 Optuna
trials is 0.064, which puts the bias-corrected floor of the 0.4301 best at \textbf{0.3662}.
And reaching an effective standard error of 0.010 needs \textbf{ten replicates per arm},
which is what sets the current run's twelve.

%% notes-tex/027-metabolic-flux-module/sections/6-preliminary.tex
%% [[experiments.027-betaxanthin-metabolic]]
%% https://github.com/Mjvolk3/torchcell/tree/main/notes-tex/027-metabolic-flux-module/sections/6-preliminary.tex

\section{The run now on Delta\secstatus{todo}}
\label{sec:preliminary}

\subsection{What was launched\secstatus{todo}}

Job array \texttt{21797666\_[0-8]} on \texttt{bbub-delta-gpu}, partition
\texttt{gpuA40x4}: nine nodes, four A40 per node, \textbf{two runs per GPU}, so 72 workers.
Six arms divide 72 into twelve seed-groups, giving twelve seeds per arm, above the ten the
noise floor requires. Each worker owns its own Optuna database, which is what avoids the
race that previously had several workers claiming one trial and silently turning four Delta
cells into duplicates.

The primary endpoint is Spearman on the pinned Merzbacher 639-gene test split at an epoch
selected on validation, so the selection set and the reporting set are disjoint. The
decision rule is the difference of differences in Section~\ref{sec:arms}.

\begin{table}[H]
  \centering
  \footnotesize
  \caption[Delta environment]{Every environment fact confirmed on the node rather than
  assumed, since the repository's own files disagree about several of them.}
  \label{tab:delta-env}
  %% SOURCE: the job's own preflight block, slurm/output/027-bxfx_21797666_0.out.
  \begin{tabular}{@{}p{0.24\textwidth}p{0.68\textwidth}@{}}
    \hline
    Fact & Value \\
    \hline
    \texttt{DATA\_ROOT} & \texttt{/scratch/bbub/mjvolk3/torchcell}. The Delta
      \texttt{.env} still says \texttt{/work/hdd}, which is wrong; pointing at a root with
      no build makes the dataset attempt a Neo4j rebuild on a compute node \\
    interpreter & \texttt{/work/hdd/bbub/miniconda3/envs/torchcell/bin/python}, called
      directly, no \texttt{conda activate} and no container \\
    GPU & NVIDIA A40, 46{,}068 MiB, four per node \\
    \texttt{NUM\_WORKERS} & 0, forced. Spawned workers re-import off Lustre and stall \\
    W\&B & online, entity \texttt{zhao-group}, project \texttt{torchcell\_027\_bxfx} \\
    \hline
  \end{tabular}
\end{table}

\subsection{Three faults found and fixed before submitting\secstatus{todo}}

\textbf{The Open Enzyme Database mirror was absent from Delta.}
\texttt{resolve\_kcat\_table} catches the resulting \texttt{FileNotFoundError} and returns
an empty record list, which replaces every measured $k_{cat}$ with the organism default. The
enzyme-constrained arm would have run with no enzyme constraints and reported a plausible
number. The mirror is 516\,kB and was synced; the coverage gate exists to abort exactly this
case, and it can only fire if there is a mirror to measure.

\textbf{\texttt{WANDB\_ENTITY} was unset}, so runs would have landed under whatever default
entity the node carried while 026's are under \texttt{zhao-group}, leaving the arms
unjoinable at analysis time.

\textbf{\texttt{OMP\_NUM\_THREADS} was unset}, so eight workers on a 32-CPU allocation of a
64-core node would each have defaulted to the node's core count.

\subsection{Three couplings the backbone change forced\secstatus{todo}}

Moving 027 onto the strong backbone broke three things, each found by running a smoke test
rather than by reading the code, and each silent in a different way.

\begin{enumerate}\setlength{\itemsep}{2pt}
  \item The perturbation head defaults to eight attention heads and 90 is not divisible by
    eight, so its cross-attention refuses to build. 019 uses six.
  \item The model reads \texttt{cell\_graph["gene"].x} only when it is told it has
    precomputed embeddings. With the learnable table disabled and \texttt{node\_embeddings}
    unset it finds neither and raises. It also indexes the embedding dataset to infer its
    width, so it needs the objects rather than their names.
  \item Dropping 026's \texttt{build\_dataset} dropped its module globals with it.
\end{enumerate}

A two-epoch local smoke run on the corrected backbone moved loss 1.0586 to 1.0088 and
validation Pearson $-0.0214$ to $+0.1055$, and computed the pinned test Spearman at
$+0.0261$ on $n = 629$. Two epochs is not a result. What it establishes is that the
nine-graph \texttt{prot\_T5} backbone builds, trains and scores end to end.

\subsection{Preliminary training results\secstatus{todo}}

\textbf{None yet.} At the time of writing, four of the nine array tasks are running and
five are queued behind them on resources. No worker has emitted a training epoch, so there
is no validation Pearson to compare and none is quoted here.

The delay is understood rather than mysterious. The launcher runs a sequential
\texttt{--create-only} pass per worker before any GPU work, both to create each worker's
Optuna database and to pay the cold-import cost once so the Lustre client cache is warm for
the rest. That first import is reading at roughly 60\,kB\,s$^{-1}$: the process sits in
\texttt{ptlrpc\_set\_wait}, the Lustre RPC wait, at about one percent CPU, having read
90\,MB after twenty minutes. It is progressing, and a 20-hour allocation absorbs an hour of
cache warming, but the observation is worth recording because the same cost has previously
been mistaken for a hang and has killed canaries with short wall limits.

This section will carry the arm comparison once the sweep produces one. The comparison that
matters is not a single validation Pearson but the difference of differences against the
permuted-label nulls, and that needs the seeds to finish.


\end{document}
